Este estudio se cierra con la versión del 13 de octubre de 2026. Las conclusiones siguientes no establecen qué ocurrió ni por qué. Describen hacia dónde apuntan los resultados de un modelo numérico, según los parámetros de entrada declarados en este documento y las 84 simulaciones completadas: 19 de cribado con el terreno anterior, 41 con la ciudad completa, 5 de cuadratura de la concentración y 19 microsimulaciones del campus. Con otros datos de entrada, otro modelo u observaciones nuevas, los resultados pueden cambiar.

\subsection*{Hacia dónde apuntan los resultados}
\begin{enumerate}
\item \textbf{Estructura general.} Con la ciudad completa y los supuestos adoptados, el modelo reproduce los corredores de Camino El Alba, General Blanche y Quebrada Honda que delimitan la zona afectada oficial, y el alcance hasta Padre Hurtado. No reproduce la distribución espacial de la afectación dentro de esa zona: la prueba con observaciones reservadas, definida antes de examinar los resultados y ejecutada una sola vez, fue desfavorable (19 de 49 equipamientos, frente al \SI{70}{\percent} fijado).
\item \textbf{Volumen.} En los escenarios ensayados, unos \SI{300000}{m^3} de mezcla permiten reproducir el alcance observado. Es el mismo orden de magnitud que una estimación independiente por balance de lluvia que adopta la saturación del suelo informada. Ninguno de los dos valores es una medición.
\item \textbf{Entradas de flujo al modelo.} Los escenarios en que el barro sale por un solo punto no reproducen el conjunto de lo observado. Los más compatibles combinan la entrada por el cauce principal, que en el modelo reproduce la rama norte, con el brazo sur y la lluvia supuesta sobre la ladera, que reproducen parte de la afectación informada en el sector sur.
\item \textbf{Papel de la ciudad.} Representar los 42.018 edificios cambia el resultado: en el modelo, las calles concentran el flujo y aumentan los espesores.
\item \textbf{Drenaje.} Los escenarios de capacidad del colector de Francisco Bulnes Correa son hipotéticos. En el modelo, con el volumen supuesto, llega a la embocadura un caudal calculado del orden de \SI{65}{m^3/s}. Con capacidades supuestas menores, el flujo continúa en superficie por Quebrada Honda. Estos escenarios no describen la obra real, para la que no se dispuso de información pública sobre capacidad, sección y estado. No permiten afirmar que la obra estuviera obstruida, que funcionara mal ni que una capacidad distinta hubiese cambiado lo ocurrido.
\item \textbf{Campus de la Universidad de los Andes.} En las simulaciones, el aporte supuesto de agua turbia de la ladera y del cerro sigue La Plaza Sur y entra por un claro a un patio hundido; ese recorrido calculado es compatible con las imágenes consultadas. Con malla fina (m36), el espesor calculado en el patio llega a \SI{0,17}{m}, en el rango estimado en el video publicado, y La Plaza Sur se moja en el \SI{58}{\percent} de sus puntos. El desborde simulado de Álvaro del Portillo no llega a ese patio con este relieve.
\item \textbf{Incendio.} Según la severidad estimada con Sentinel-2, el \SI{82}{\percent} de las cuencas aportantes muestra señal de quema y el \SI{10}{\percent} tiene severidad moderada-alta o alta. Con los coeficientes supuestos, el balance sin incendio da alrededor del \SI{62}{\percent} del volumen de barro y el \SI{43}{\percent} de los sólidos del escenario con incendio. Staley M1, de uso ilustrativo, da una probabilidad mediana de 0,28 con incendio frente a 0,07 sin él. En la ciudad, las repeticiones corregidas sin incendio (f35b y f38b) mojan \SIrange{179}{193}{ha}, frente a \SI{222}{ha} en p34b (\SIrange{13}{19}{\percent} menos), y calculan \SIrange{0,19}{0,23}{m} junto a Clínica UC Christus. El IEC queda en 80--82, frente a 81 con incendio, y no distingue esos escenarios.
\item \textbf{Mejor ajuste.} La hipótesis integrada obtiene un índice exploratorio de compatibilidad (IEC) de 79--82/100: 82 con \SI{35}{mm/h} supuestos y paso fino (z32), 81 con el pulso medido en Río Molina (p34b, corregida) y 79 con esa lluvia repartida en \SI{4}{h} (v33). Combina cauces repartidos, lluvia en el faldeo, casas que dejan entrar el barro, agua turbia en el faldeo y los rasgos del campus. Es un ajuste posterior a la prueba reservada: no está validado y el índice no es una medida de precisión.
\item \textbf{Hipótesis sin respaldo en el modelo.} En los ensayos, tallar la pista del \emph{bike park} no cambia el recorrido calculado del agua, y un obstáculo supuesto de troncos al final de la línea de árboles no desvía el flujo simulado hacia el sur.
\end{enumerate}

\subsection*{Lo que permanece sin probar}
\begin{itemize}
\item \textbf{Validación independiente.} Ninguna hipótesis posterior a la prueba reservada se contrastó con observaciones nuevas, por ejemplo radar Sentinel-1, marcas de nivel o espesores medidos en terreno.
\item \textbf{Lluvia en la cuenca.} No se identificó un pluviómetro dentro de la cuenca en las fuentes consultadas. Se usó la estación Río Molina, a unos \SI{7}{km}, con un desfase temporal supuesto.
\item \textbf{Volumen y composición reales del flujo.} No se localizaron mediciones del volumen movilizado, de su concentración ni de su reología en las fuentes consultadas. Permanecen no probadas las distribuciones alternativas a la Beta(3,2) usada en la cuadratura y la estimación de probabilidades de inundación y percentiles extremos.
\item \textbf{Obras de drenaje.} Se desconocen la capacidad, la sección y el estado reales del colector y de los ductos. Los escenarios del estudio no los evalúan.
\item \textbf{Clínica UC Christus.} Ninguna configuración reproduce a la vez el espesor cercano a \SI{1}{m} informado y la extensión dentro de la zona afectada oficial. La calibración de la retención en viviendas con información de daños por vivienda permanece no probada.
\item \textbf{Relieve.} El modelo de elevación de \SI{30}{m} no resuelve calzadas, soleras ni pasos estrechos. Los rasgos tallados son representaciones gruesas basadas en observaciones aportadas. No se dispuso de un relieve LiDAR.
\item \textbf{Procesos no simulados:} erosión y depósito durante el recorrido, troncos y vehículos arrastrados, drenaje urbano, infiltración dentro de la ciudad y cambios de concentración del flujo. Permanecen no probados el trazador que separaría el aporte de los hidrogramas de la lluvia, el sedigrama y el transporte de concentración.
\item \textbf{Parámetros del incendio.} Los incrementos de escorrentía por severidad son supuestos; su calibración local y el tiempo de concentración de cada subcuenca según severidad permanecen no probados. Los umbrales de severidad y los modelos empíricos de volumen y probabilidad se calibraron fuera de Chile.
\item \textbf{Obstáculo hipotético de árboles.} Solo se ensayó una ubicación y una altura supuestas, sobre el relieve disponible. No se verificó en terreno.
\end{itemize}

\subsection*{Naturaleza y alcance de este estudio}
Este documento es un ejercicio técnico independiente, de carácter académico y divulgativo, realizado por iniciativa personal y sin encargo de terceros. No es un peritaje, un informe pericial, una opinión técnica oficial ni una evaluación de obras, servicios o decisiones de ninguna persona o institución. Se basa exclusivamente en información pública consultada entre el 8 y el 13 de octubre de 2026, registrada con fecha y huella en el registro probatorio. El autor no tiene conocimiento directo de los hechos más allá de esas fuentes.

Los resultados son respuestas de un modelo numérico a parámetros de entrada declarados y a supuestos explícitos, varios de ellos hipotéticos. Ninguna cifra calculada es una medición, y ninguna frase de este documento afirma la causa del evento ni atribuye responsabilidades. Las menciones a obras, servicios o instituciones se refieren a escenarios del modelo o a lo que informaron las fuentes citadas, con su atribución. Las fuentes discrepan en varios puntos, y esas discrepancias se presentan sin resolverlas.

El estudio queda cerrado en esta versión. El código, los datos derivados y las versiones congeladas están disponibles para su revisión y reproducción. Cualquier uso de estos resultados en decisiones técnicas, administrativas o judiciales requiere una evaluación profesional independiente, con información de diseño, inspección y mediciones en terreno.
